\clearpage
\section{Experimentalist: testing the methods}
\label{sec:experiment}

\subsection{Benchmark design}

Real sensor archives rarely provide a certain answer about every reading. To
compare methods directly, I built a synthetic month of air temperature with
30 days of five-minute samples, giving 8640 readings. The clean signal
contains a daily cycle, a slower weather component and correlated sensor
noise. I then inserted five fault types based on the fault categories
discussed by \citet{campbell2013} and \citet{leigh2019}:

\begin{itemize}
  \item 14 sudden spikes between 3 and 8 K;
  \item a $+2.5$ K level shift lasting 36 hours;
  \item a frozen sensor for 6 hours with exact repetition;
  \item a frozen sensor for 6 hours with $\pm0.02$ K jitter;
  \item a slow calibration drift reaching $+2.2$ K over four days.
\end{itemize}

Three healthy neighbouring stations share the same broad weather pattern but
have their own offsets and noise. Random-number generators are seeded, so the
benchmark and figures are exactly reproducible.

\begin{figure}[h]
  \centering
  \includegraphics[width=0.95\textwidth]{fig_series.png}
  \caption{Synthetic temperature benchmark. Shaded regions are segment
  faults; circles mark the injected spikes.}
  \label{fig:series}
\end{figure}

For spikes, I count a fault as detected if a flag falls within one sample of
the injected spike. For segment faults, I report the percentage of faulty
readings flagged. False alarms are flags on clean readings and are reported
per day for one series.

\subsection{Masking: why a classical test can miss two outliers}

I first reproduced the masking problem. Eighteen values were placed tightly
around 10.0 and two equal outliers at 10.6 were added. The pair increases the
sample standard deviation enough that the largest Grubbs statistic is about
2.56, below the critical value of about 2.71 for this sample. An iterated
single-outlier Grubbs procedure therefore stops and finds nothing. Rosner's
generalised ESD, which evaluates several possible removals before deciding,
finds the two outliers \citep{rosner1983}.

\begin{figure}[h]
  \centering
  \includegraphics[width=0.92\textwidth]{fig_grubbs.png}
  \caption{Left: Grubbs critical values checked by simulation. Right: the
  masking example, where a pair defeats the iterated single-outlier test but
  generalised ESD finds both.}
  \label{fig:grubbs}
\end{figure}

This experiment is useful because the failure comes from the reference
statistics, not from complicated data. Two bad values can change the
standard deviation used to judge themselves.

\subsection{The modified z-score: a good statistic with a context problem}

On the same small contaminated sample, the robust modified z-score separates
the two planted outliers because the median and MAD barely move. In that
setting, the method behaves as intended.

The problem appears when the same score is used in a short moving window over
a correlated sensor stream. I tested windows of 7, 25 and 97 samples and
swept the threshold from 2 to 400. At a 7-sample window, the textbook value
3.5 detects all injected spikes but produces 1.63 false alarms per day for a
single series. A threshold of 30 removes those false alarms but catches only
9 of the 14 spikes.

\begin{figure}[h]
  \centering
  \includegraphics[width=0.92\textwidth]{fig_thresholds.png}
  \caption{False alarms and spike detection as the modified z-score
  threshold changes. The dashed line is 3.5. No threshold is both silent and
  complete for the tested windows.}
  \label{fig:thresholds}
\end{figure}

The important point is not that 3.5 is ``wrong'' or 30 is ``right''. The
point is that a threshold calibrated for one data setting should not be
copied into another without measuring what it costs. Window length, sampling
rate, noise and threshold form one design problem.

\subsection{Comparing the checks}

I then ran the main implementable checks over the same benchmark. Internal
consistency is not scored because the benchmark has only one measured
quantity per station. The feature-based branch from \citet{talagala2019} is
also not implemented because fourteen events are too few for a meaningful
evaluation of its extreme-value machinery.

Table~\ref{tab:results} gives the exact comparison. The final row is a simple
battery that takes the union of the step, variance-persistence and spatial
checks.

\begin{table}[h]
\centering
\small
\setlength{\tabcolsep}{2.8pt}
\caption{Percentage of faulty readings flagged, with false alarms per day.}
\label{tab:results}
\begin{tabular}{@{}lrrrrrr@{}}
\toprule
Check & Spike & Shift & Exact & Jitter & Drift & False/day \\
\midrule
Range & 0.0 & 0.0 & 0.0 & 0.0 & 0.0 & 0.00 \\
Grubbs, whole month & 0.0 & 0.0 & 0.0 & 0.0 & 0.0 & 0.00 \\
Modified z, whole month, 3.5 & 0.0 & 0.0 & 0.0 & 0.0 & 0.0 & 0.00 \\
Modified z, win. 7, 3.5 & 100.0 & 0.9 & 0.0 & 2.8 & 1.0 & 1.63 \\
Modified z, win. 7, 30 & 64.3 & 0.2 & 0.0 & 0.0 & 0.0 & 0.00 \\
Step, 3 K/5 min & 100.0 & 0.5 & 0.0 & 0.0 & 0.0 & 0.50 \\
Persistence, exact & 0.0 & 0.0 & 100.0 & 0.0 & 0.0 & 0.00 \\
Persistence, variance & 0.0 & 0.0 & 100.0 & 100.0 & 0.0 & 0.00 \\
Spatial, neighbours & 100.0 & 99.8 & 80.6 & 81.9 & 32.4 & 0.10 \\
Forecast, one-step & 100.0 & 0.7 & 0.0 & 0.0 & 0.0 & 0.57 \\
\textbf{Battery: step + persist. + spatial} & \textbf{100.0} & \textbf{100.0} & \textbf{100.0} & \textbf{100.0} & \textbf{32.4} & \textbf{0.60} \\
\bottomrule
\end{tabular}
\end{table}

\begin{figure}[h]
  \centering
  \includegraphics[width=0.95\textwidth]{fig_detectors.png}
  \caption{The same comparison as a detection matrix. No single check covers
  all fault types; the battery covers the first four completely and retains
  the spatial check's partial sensitivity to drift.}
  \label{fig:detectors}
\end{figure}

The results are easier to understand by fault type than by method. Spikes
are easy for step, rolling outlier and forecast checks. Exact persistence is
perfect for a perfectly frozen sensor but fails completely when the stuck
value jitters. The low-variability persistence check fixes that problem.
Spatial comparison is the strongest single detector because healthy
neighbours provide an external reference. Slow drift remains difficult and
is detected only after the accumulated difference becomes large enough.

The battery is the practical result I take most seriously. Three simple and
explainable checks cover more faults than any individual method. This agrees
with the operational philosophy of \citet{shafer2000} and the combination
recommended by \citet{leigh2019}.

\subsection{Limits of the experiment}

The benchmark is deliberately cleaner than the real world. It contains no
missing data, fouling that resembles natural behaviour, or local storms that
make one healthy station disagree with its neighbours. The neighbouring
stations are also more cooperative than real stations separated by complex
terrain. Therefore the exact percentages are properties of this benchmark,
not universal performance claims.

The constants are also chosen deliberately rather than derived: a 3 K step
limit, a 0.02 K variability floor, a three-standard-deviation spatial band,
and a wider forecast band. This is not hidden because it is part of the
point. A threshold is an engineering choice unless its assumptions and cost
have been demonstrated for the deployment.


\subsection{Reading the results by fault type}

The result table is easier to remember when it is read by fault rather than
by algorithm. This also matches the practical question an operator has:
``if this kind of sensor problem happens, which evidence is likely to find
it?''

\paragraph{Spikes.}
Spikes are the easiest fault in the benchmark. The step check catches all of
them because a spike creates an abrupt change. The rolling modified z-score
at threshold 3.5 and the one-step forecast also catch all of them. Spatial
comparison catches them because the neighbours remain healthy. This is a
good example of several independent checks agreeing on the same event.
Agreement does not make the value mathematically impossible, but it gives a
reviewer much stronger evidence than one flag alone.

\paragraph{Level shift.}
The planted level shift is $+2.5$\,K while the step limit is 3\,K per five
minutes. For that reason the step check is almost blind to it. The forecast
check notices mainly the entry to the new level and then adapts. Spatial
comparison performs much better because the target station stays separated
from the healthy neighbours for the duration of the fault. This illustrates
the value of an external reference: the station's own recent history may
accept a new level that its neighbours show to be suspicious.

\paragraph{Frozen sensor with exact repetition.}
Exact persistence is designed for this case and detects it completely. The
variance form also detects it because a perfectly repeated value has zero
spread. Spatial comparison detects most of the period once the real weather
moves away from the frozen value. The result is straightforward: when a
fault matches the definition of the check, even a very simple rule can be
excellent.

\paragraph{Frozen sensor with small jitter.}
This fault is deliberately almost the same as the previous one, except that
the last digit moves by $\pm0.02$\,K. That tiny change is enough to defeat
exact equality. The low-variability persistence check still detects the
whole fault, because the window is much less variable than a healthy sensor.
This is one of the clearest lessons in the report: a small implementation
detail can decide whether an apparently standard check works or fails.

\paragraph{Slow drift.}
Drift is the difficult column. Each individual change is small, so the step
check sees nothing. A moving outlier reference and a one-step forecast move
with the fault, which also makes them weak. Spatial comparison is the only
implemented check with meaningful sensitivity, and even it detects the drift
late, after enough disagreement has accumulated. The correct conclusion is
not that spatial checking solves drift; it is that drift needs a reference
that does not drift with the faulty sensor.

\paragraph{False alarms.}
Detection is only half of the result. A check that catches every fault but
creates many false alarms may be unusable in a large network. At the
textbook modified-z threshold of 3.5, the rolling detector produces 1.63
false alarms per day for one series in this benchmark. Multiplying that rate
by many sensors quickly creates reviewer workload. The battery produces 0.60
false alarms per day here while covering the first four fault classes
completely. The exact number is benchmark-specific, but the trade-off is the
important part.

\subsection{Why the battery is the main practical result}

The battery combines step, variance-persistence and spatial checks. I chose
these three because their strengths are different and their rules are easy
to explain. Step handles abrupt movement, variance-persistence handles a
sensor that stops moving, and spatial comparison supplies an external
reference. Their union therefore covers faults that any one of them misses.

This does not prove that these three checks are optimal for every network.
It shows something more modest and more useful: combining transparent checks
with different failure modes can outperform a more ambitious search for one
universal anomaly detector. The same principle appears in the operational
literature, where several checks feed a common flagging and review process.

In a production system I would therefore evaluate a battery in the same way
as an individual detector. I would measure detection by fault type, false
alarms per reviewer per day, the number of checks that agree on each event,
and the cases that still escape all checks. This keeps the evaluation tied
to the operational decision rather than to one abstract accuracy number.


\subsection{Parameters used in the benchmark}

To keep the experiment explainable, the benchmark uses a small number of
explicit parameters. Table~\ref{tab:simpleparams} collects the ones that
matter most. These values are not presented as universal recommendations;
they are the operating point used to compare the methods on the same data.

\begin{table}[h]
\centering
\small
\caption{Main benchmark and detector parameters used in the experiment.}
\label{tab:simpleparams}
\begin{tabular}{@{}p{4.4cm}p{3.6cm}p{6.0cm}@{}}
\toprule
Item & Value & Reason in this benchmark \\
\midrule
Sampling interval & 5 minutes & Matches the high-frequency sensor setting used in the experiment. \\
Temperature range & $-40$ to $50$\,\textcelsius & Broad physically plausible range; intended only to catch gross errors. \\
Step limit & 3\,K per 5 minutes & Large enough to ignore ordinary movement but small enough to catch the planted spikes. \\
Exact persistence & 12 repeated samples & One hour of identical readings. \\
Variance persistence & 3-hour window; 0.02\,K spread floor & Detects a sensor that is almost constant without requiring exact equality. \\
Spatial reference & 3 healthy neighbours & Supplies an external reference independent of the faulty target sensor. \\
Spatial band & 3 standard deviations & Simple conventional operating point for the comparison. \\
Modified z-score & windows 7, 25, 97 & Shows how the same score changes when its reference window changes. \\
Forecast training & first 5 days & Gives the simple forecaster repeated examples of the daily pattern. \\
\bottomrule
\end{tabular}
\end{table}

There are two reasons for writing these values down. First, a result cannot
be reproduced if the threshold and window are hidden. Second, the meaning of
a threshold depends on its context. A 3\,K step limit is a statement about
a five-minute sampling interval; it would mean something different for
hourly data. A persistence window that is reasonable for air temperature may
be far too short for deep soil temperature. Even a standard-deviation band
changes meaning when the noise and the training period change.

For this reason I treat the parameter table as part of the method rather
than as implementation detail. If the system were deployed at another site,
I would keep the checks but re-evaluate the values using validated local
data, instrument specifications and an acceptable false-alarm budget. This
is also why the experiment reports the parameters beside the results instead
of presenting the percentages as if they came from the algorithms alone.

% =====================================================================
% 5. Futurist
% =====================================================================
